% Statement and cited-input packet. Target proof is separately embargoed.
% Frozen II source SHA256: 74497bfe327fae25c9cdcb17334b3bf42a9793a9df7c0a4404d97568f2f760f0
% Target owner: push:src:E.2.13; tag: 032Y

\begin{proposition}[persistence of an isolated finite multiple fibre]\hypertarget{tag-032Y}{}\label{push:src:E.2.13}


In the setting of \ref{push:src:E.2.10}, suppose the central value of
one finite labelled critical section is distinct from
the limits of the other two labelled sections.  Write
its multiplicity as $m$, with $m=3k$ or $4l$, and its
linear monodromy order as $e=3$ or $4$, respectively.
Near that section, normalize the base change taking
an $m$th root of its defining equation.

The resulting total space is smooth, and its map to
the root bidisc is a proper smooth family of two-tori.
The deck group acts freely on the total space.
Consequently the actual central reduced finite fibre
is smooth, has exactly multiplicity $m$, and belongs
to the same smooth family of finite free affine
quotients as the punctured reduced finite fibres.

This applies to both finite fibres when the three
labelled critical values have distinct limits.  In
the two-critical-value collapsed-target case, it
applies to the isolated finite labelled fibre.
\end{proposition}

% Cited input push:src:E.2.10; statement or labelled formula only.
\begin{proposition}[persistence for either canonical slope]\hypertarget{tag-032V}{}\label{push:src:E.2.10}


Under the hypotheses of \ref{push:src:E.2.5}, for every nonzero
$D=kl-3k-4l$ one has $E_0=0$ and $e=1$.
After shrinking the disc, the evaluation of
$\pi_*\mathcal L$ gives a proper flat holomorphic map
\[
 g:\mathcal X\longrightarrow\mathbb P^1\times\Delta
\]
whose restriction to every fibre is its intrinsic algebraic
reduction.  A smooth general central fibre of $g$ is a
complex two-torus.  Its normalized punctured period has
an interior limit as in \ref{push:src:E.2.7}.

Two families satisfying \ref{push:src:E.2.5} cannot have the same smooth
central manifold unless their ordered multiplicity pairs
agree, regardless of the signs of their canonical slopes.
Uniqueness of a smooth completion at an interior period and
the arbitrary deformation-chain statement \ref{push:src:E.2.4} are not
asserted.
\end{proposition}

% Cited input push:src:E.2.11; statement or labelled formula only.
\begin{proposition}[purity with a smooth torus Jacobian]\hypertarget{tag-032W}{}\label{push:src:E.2.11}


Let $B\subset\C^2$ be a sufficiently small bidisc about $0$, put
$U=B\setminus\{0\}$, and let $g:X\to B$ be a proper flat
holomorphic map from a connected smooth complex $(n+2)$-fold,
with connected fibres of dimension $n$.  Let $G\to B$ be a
smooth proper holomorphic family of complex $n$-tori with a
holomorphic zero section and its group law.  Suppose that
$X_U\to U$ is a holomorphic $G_U$-torsor.  Then $g$ is a
submersion on all of $X$, and the given action extends uniquely
to make $X\to B$ a holomorphic $G$-torsor.

No polarization, K\"ahler metric on $X$, or section of the
given torsor on $U$ is assumed.  The smooth extension $G$ is
an explicit hypothesis.
\end{proposition}

% Cited input push:src:E.2.12; statement or labelled formula only.
\begin{proposition}[no additional central critical values]\hypertarget{tag-032X}{}\label{push:src:E.2.12}


In the smooth-specialization setting of \ref{push:src:E.2.10}, let
\[
 g:\mathcal X\longrightarrow\mathbb P^1\times\Delta
\]
be the extended intrinsic fibration, and let
$c_1(s),c_2(s),c_0(s)$ be the three labelled critical-value
sections of \ref{push:src:E.2.7}.  The critical values of $g_0$ are exactly
the distinct points among $c_1(0),c_2(0),c_0(0)$.
In particular, no additional critical value occurs only
on the central target curve.

If the normalizing projective transformations $A_s$ of
\ref{push:src:E.2.7} have a degenerate limit, the central regular fibres
are isomorphic complex two-tori and $g_0$ has at most two
critical values.  If the constant image of the limiting
rank-one matrix is $0$ or $1$, this torus is the
corresponding finite-root abelian surface.  If that image
is an ordinary source point, it is the ordinary torus at
that point and need not be an abelian surface.  The cusp
cannot be the constant image.
\end{proposition}

% Cited input push:src:E.2.4; statement or labelled formula only.
\begin{corollary}[arbitrary connected deformation chains]\hypertarget{tag-033C}{}\label{push:src:E.2.4}


The ordered coprime pair $(3k,4l)$ is constant in every
proper smooth holomorphic family over a connected complex
analytic base which contains a constructed sphere member.
Two constructed members are deformation equivalent through
such families, allowing arbitrary intermediate members and
finite chains, if and only if their pairs agree.
Moreover different pairs cannot occur in one continuous
family in the $C^\infty$ space of integrable complex
structures on $S^6$.
\end{corollary}

% Cited input push:src:E.2.5; statement or labelled formula only.
\begin{proposition}[a fractional canonical pencil at a smooth limit]\hypertarget{tag-033D}{}\label{push:src:E.2.5}


Let $\pi:\mathcal X\to\Delta$ be a proper holomorphic submersion
whose fibres are diffeomorphic to $S^6$.  Suppose that for every
$t\ne0$ the fibre is one of the members of \SeriesRef{push:src:E.2.1}, with the same
coprime multiplicities $(3k,4l)$.  Set
\[
 A=12kl,\qquad D=kl-3k-4l\ne0.
\]
There is a unique holomorphic line bundle $\mathcal L$ on
$\mathcal X$, up to isomorphism, satisfying
\[
 \mathcal L^{D}=K_{\mathcal X/\Delta}^{A}.
\]
Its restriction to $X_t$ for $t\ne0$ is the coarse pencil
$f_t^*\mathcal O_{\mathbb P^1}(1)$.  The coherent sheaf
$\pi_*\mathcal L$ is locally free of rank two, and its central
fibre injects into $H^0(X_0,\mathcal L_0)$.

The central manifold has algebraic dimension one and a
holomorphic algebraic reduction $g_0:X_0\to\mathbb P^1$ with
connected fibres.  If $s_0,s_1$ are the two central sections
obtained from a local frame of $\pi_*\mathcal L$, there are a
nonconstant rational map $R:\mathbb P^1\to\mathbb P^1$ of
degree $e\ge1$ and an effective Cartier divisor $E_0$ such that
\[
 [s_0:s_1]=R\circ g_0,\qquad
 \mathcal L_0\simeq g_0^*\mathcal O(e)\otimes\mathcal O(E_0).
\]
Here $[s_0:s_1]$ denotes the holomorphic extension of the ratio;
the displayed divisor is its common fixed divisor.  The
statement does not assert that $E_0=0$.
\end{proposition}

% Cited input push:src:E.2.7; statement or labelled formula only.
\begin{proposition}[an extended fibration cannot reach the zero period]\hypertarget{tag-033F}{}\label{push:src:E.2.7}


Under the hypotheses of \ref{push:src:E.2.5}, suppose in addition that the
intrinsic algebraic reductions extend to a holomorphic map
\[
 g:\mathcal X\longrightarrow\mathbb P^1\times\Delta
\]
whose restriction to $X_0$ is nonconstant.  This hypothesis
holds in \ref{push:src:E.2.6}.  Normalize the punctured family to fixed
discrete logarithmic data with multiplicities $(3k,4l)$, and
write its isomorphism parameter as
\[
 z=\exp\!\left(-\frac{\pi i c}{kl}\right).
\]
The resulting holomorphic function $z:\Delta^*\to\mathbb C^*$
extends across zero and satisfies $z(0)\ne0$.  Its central
value lies strictly inside the admissible period disc.

This statement concerns the period function and the extended
relative fibration.  Identification of $X_0$ with the member
at $z(0)$ is proved separately in \ref{push:src:E.2.20}--\ref{push:src:E.2.21}; propagation
to arbitrary deformation bases and continuous families is
proved in \ref{push:src:E.2.22}--\ref{push:src:E.2.24}.
\end{proposition}

% Cited input push:src:E.2.20; statement or labelled formula only.
\begin{proposition}[identification of the actual analytic cusp]\hypertarget{tag-0336}{}\label{push:src:E.2.20}


In the smooth specialization of \ref{push:src:E.2.10}, normalize
the three distinct critical-value sections to
$0,1,\infty$ and write $q=1/t$ near the cusp.
Let $c(s)$ be the holomorphic interior period
lift supplied by \ref{push:src:E.2.7}.  After shrinking the
$(q,s)$-bidisc, its inverse image in the actual
family is biholomorphic, over that bidisc, to
the prescribed marked Mumford cusp family at
parameter $c(s)$.  The isomorphism induces
the given identification of the regular
torus Jacobians and their integral markings.
Its affine origin may be fixed by choosing
a holomorphic section through a smooth
point of the reduced cusp.
\end{proposition}

% Cited input push:src:E.2.21; statement or labelled formula only.
\begin{proposition}[global identification of a smooth completion]\hypertarget{tag-0337}{}\label{push:src:E.2.21}


Let $\pi:\mathcal X\to\Delta$ satisfy the hypotheses of
\ref{push:src:E.2.5}, with fixed coprime multiplicities $(3k,4l)$ on
every noncentral fibre.  Then $X_0$ is a member of that
same constructed family.  Its isomorphism parameter is
the interior value $z(0)$ of \ref{push:src:E.2.7}.

After choosing a holomorphic lift $c(s)$ of that parameter
and shrinking the disc, the actual family is isomorphic
to the prescribed family at $c(s)$.  The ambiguity in
this comparison is the connected translation automorphism
group already determined in E.1.
\end{proposition}

% Cited input push:src:E.2.22; statement or labelled formula only.
\begin{proposition}[propagation from a known open subset of a disc]\hypertarget{tag-0338}{}\label{push:src:E.2.22}


Let $\pi:\mathcal X\to\Delta$ be a proper holomorphic
submersion with fibres diffeomorphic to $S^6$.  Suppose
there is a nonempty open subset $U\subset\Delta$ such
that every $X_s$, $s\in U$, is a member of the fixed
coprime multiplicity family $(3k,4l)$ of \SeriesRef{push:src:E.2.1}.
Then every fibre over $\Delta$ is a member of that
same family.
\end{proposition}

% Cited input push:src:E.2.24; statement or labelled formula only.
\begin{corollary}[separation for continuous complex structures]\hypertarget{tag-033A}{}\label{push:src:E.2.24}


Let $\mathcal I(S^6)$ be the space of smooth integrable
almost-complex structures on the fixed smooth sphere,
with its $C^\infty$ topology.  For a fixed coprime pair
$(3k,4l)$, let $\mathcal I_{k,l}$ consist of the structures
whose complex manifolds are biholomorphic to constructed
members with that pair.  Then $\mathcal I_{k,l}$ is open
and closed in $\mathcal I(S^6)$.

Consequently a continuous family of such structures
over any connected topological parameter space cannot
contain constructed members with two different pairs.
The same separation holds in the quotient spaces by
$\operatorname{Diff}_0(S^6)$ or by all orientation-preserving
diffeomorphisms, with their quotient topologies.
This statement does not identify the number of marked
connected components belonging to a single pair.
\end{corollary}

% Cited input push:src:E.2.6; statement or labelled formula only.
\begin{proposition}[positive canonical slope preserves the fibration]\hypertarget{tag-033E}{}\label{push:src:E.2.6}


Under the hypotheses of \ref{push:src:E.2.5}, suppose $D>0$.  Then $E_0=0$
and $e=1$.  After shrinking the disc and choosing a frame of
$\pi_*\mathcal L$, its evaluation defines a proper flat
holomorphic map
\[
 g:\mathcal X\longrightarrow\mathbb P^1\times\Delta
\]
which is the intrinsic algebraic reduction on every fibre.
Its smooth general central fibre is a complex two-torus.

In particular, two families of the type in \ref{push:src:E.2.5} with positive
canonical slope cannot have the same smooth central manifold
unless their multiplicity pairs agree.  A positive-slope and
a negative-slope family cannot have a common smooth central
manifold either.  This proposition concerns a common smooth
specialization.  The additional propagation and closure
arguments for arbitrary connected chains are \ref{push:src:E.2.22}--\ref{push:src:E.2.24}.
\end{proposition}

